%\chapter{ABSTRACT}
\clearpage
\phantomsection
\addcontentsline{toc}{chapter}{ABSTRACT}
\begin{center}
\textbf{\large  ABSTRACT}\\[0.5cm]
\end{center}

This thesis presents a comprehensive investigation into the aerothermal characteristics of Hypersonic Inflatable Aerodynamic Decelerators (HIADs), with a particular focus on the impact of their scalloped surface topology on flow physics and aerodynamic heating. Utilizing Computational Fluid Dynamics (CFD), the study reviews related research in HIAD development and re-entry aerothermodynamics. It delves into the fundamental principles of hypersonic aerodynamics, including shock wave formation, viscous interactions, and critical heat transfer mechanisms. The work also compares conventional Thermal Protection Systems (TPS) with flexible HIAD structures, highlighting the benefits and challenges of the latter. Furthermore, the theoretical framework of CFD, turbulence modeling (specifically the SST k-omega model), and advanced thermodynamic gas models are discussed to establish a robust methodology for analyzing complex hypersonic flow fields and predicting localized peak heating effects.




\noindent \textbf{Keywords}: Hypersonic Inflatable Aerodynamic Decelerator (HIAD), Aerothermodynamics, Computational Fluid Dynamics (CFD), Shock-Boundary Layer Interaction (SBLI), Scalloped Geometry, Aerodynamic Heating, Thermally Perfect Gas